\section{Introduction}
\label{sec:introduction}

Quantum computing processes information through superposition, interference, and entanglement, offering alternatives to classical computation for quantum simulation and selected computationally intensive problems~\cite{preskill2018quantum}. Executing larger circuits requires more qubits, more precise operations, and hardware control and compilation that limit accumulated execution errors. Neutral-atom processors confine atoms in optical-tweezer arrays, implement entangling gates through Rydberg interactions, and reconfigure connectivity through coherent transport. Experimental advances in arrays with hundreds of qubits, high-fidelity entangling gates, and coherent atom transport~\cite{bluvstein2024logicalprocessor,evered2023highfidelity,bluvstein2022coherenttransport} make neutral atoms an important platform for large-scale quantum computing.

Neutral-atom compilers translate logical circuits into physical operations by arranging atom positions, trajectories, and laser pulses. Zoned neutral-atom architectures separate atom storage from entangling-gate execution, using storage isolation to suppress Rydberg excitation of idle atoms. Inter-zone transport incurs trap-transfer losses and consumes coherence time, making inter-layer placement and transport scheduling important to circuit fidelity and execution time~\cite{stade2024abstractmodel,lin2025zac}.

Layer-boundary atom positions affect the execution of subsequent gate layers. Retaining a temporarily idle atom in the entanglement zone avoids round-trip trap transfers through storage. The atom also encounters Rydberg pulses in intermediate layers and incurs additional excitation. Returning an atom to storage sets its later re-entry origin and may obstruct other atoms' transport paths. Structural constraints on mobile traps determine trajectory compatibility, and the longest trajectory determines each rearrangement batch's duration~\cite{lin2025zac,stade2025routingaware}. Gate sites, residency, and storage sites jointly affect physical costs in the current and subsequent layers.

ZAC, proposed by Lin et al., uses circuit-wide interactions for initial placement~\cite{lin2025zac}. Its dynamic gate-site and storage-site selection uses adjacent-layer reuse matching and next-partner distances. Stade et al. incorporate trajectory compatibility and batch duration into the A* cost for routing-aware placement~\cite{stade2025routingaware}. Evaluating residency between nonadjacent gate layers requires tracking atom positions and pulse exposure throughout the intermediate layers. These states determine how current residency and return decisions affect subsequent excitation losses, trap transfers, and feasible trajectories.

We propose GA-LK for joint layer-boundary optimization\footnote{The implementation and experimental configurations are available in the project repository at \url{https://github.com/zhouzixiang1/zac-ga}.}. Genetic search jointly selects gate sites and residency states, while injective matching assigns storage sites to returning atoms. Trajectory conflict graphs define rearrangement batches, whose execution order is checked. Decayed multilayer physical look-ahead evaluates subsequent transport, idle-atom excitation, and coherence losses from each candidate's resulting state. The search coordinates current placement and nonadjacent interaction requirements under a shared physical objective.

We evaluate GA-LK on the 18 circuits used in the ZAC paper and 154 input circuits from the MQT QMAP examples. GA-LK achieves higher geometric-mean model fidelity on both benchmark sets than Lin et al.'s ZAC compiler~\cite{lin2025zac} and Stade et al.'s routing-aware placement~\cite{stade2025routingaware}.
